%! Author = denis
%! Date = 28.06.2026

% Preamble
\documentclass[11pt]{article}

% Packages
\usepackage{amsmath}
\usepackage[margin=2.5cm]{geometry}
\usepackage{booktabs}
\usepackage{array}

\title{Cost and Time Estimation\\\large Sensor Data Collector (Mobile App)}
\author{Denis Schatzmann}
\date{\today}

% Document
\begin{document}

\maketitle

\section{Task}
The Sensor Data Collector lets a user pick sensors from a select box, add them
to a list (duplicates ignored), remove a marked sensor, and toggle data
collection with a Start/Stop button. On stop the collected data is sent to a
backend service. This report follows the steps from the lecture: estimate the
size, derive effort and schedule with COCOMO, and finally the cost.

\section{Size}
This is a small app that is mostly user interface with very little logic in the
background. For that case the lecture recommends the \emph{widget point} method,
so we use it as the main estimate and the function point count as a cross-check.

We count the interface elements:

\begin{table}[h]
\centering
\begin{tabular}{l c}
\toprule
\textbf{Widget type} & \textbf{Count} \\
\midrule
Input widgets (select box, add, remove, Start/Stop)  & 4 \\
Composite widgets (sensor list)                      & 1 \\
Describing widgets (labels)                          & 3 \\
\midrule
\textbf{Widget points} & \textbf{8} \\
\bottomrule
\end{tabular}
\caption{Widget point count.}
\end{table}

\noindent
With $\text{FP} = 2 \times \text{widget points}$ this gives \textbf{16 function
points}. A classic function point count over the same six functions (the three
buttons as external inputs, the backend send as an external output, the select
box as an external inquiry, and the sensor list as an internal file) comes out
somewhat higher, at about 23 FP. Since the app is essentially a thin user
interface over an existing backend, we take the lower widget-point value of
16 FP as the estimate and treat the 23 FP as a conservative upper bound.

\section{Effort: COCOMO}
Converting to lines of code with the language table (53 LOC/FP for Java/Kotlin)
gives $16 \times 53 \approx 850$ LOC, i.e.\ about $0.85$ KDL. Applying COCOMO in
organic mode results in a base effort of roughly \textbf{2.7 person months}.

Because this is medical software, the required reliability is rated high, which
adds about 15\,\%. The adjusted effort is therefore about \textbf{3.1 person
months}. Split over the usual phases of a small project:

\begin{table}[h]
\centering
\begin{tabular}{l c c}
\toprule
\textbf{Phase} & \textbf{Share} & \textbf{Person months} \\
\midrule
Design                & 19\% & 0.6 \\
Implementation        & 63\% & 2.0 \\
Integration and tests & 18\% & 0.6 \\
\bottomrule
\end{tabular}
\caption{Effort distribution for a small project.}
\end{table}

\section{Schedule}
Using the organic schedule formula, 3.1 person months stretch over about
\textbf{3.8 months} with an average staffing of roughly one person. So a single
developer working a bit under four months covers the work.

\section{Cost}
With a monthly salary of CHF 10'000, the personnel cost is
\[
3.1 \times 10\,000 \approx \text{CHF } 31\,000.
\]
Adding about CHF 2'000 for tooling licences, test devices and the cloud backend
gives a total of about \textbf{CHF 33'000} for one developer over roughly
\textbf{3.8 months}.

\section{Accuracy}
At this early stage the estimate is only accurate within a factor of about
$0.5$ to $2$, so the realistic range is roughly CHF 17'000 to CHF 66'000. The
estimate assumes one cross-platform codebase, an existing backend API, and
standard sensor access without custom drivers.

\end{document}
